\documentclass[11pt]{article}
\usepackage[margin=1in]{geometry}
\usepackage{booktabs}
\usepackage{amsmath}
\usepackage{parskip}
\usepackage{enumitem}

\title{Issue \#54: Embedding Projections $\rightarrow$ Facial Emotion\\[6pt]
\large Summary of Specifications, Results, and Interpretation}
\author{}
\date{April 9, 2026}

\begin{document}
\maketitle

\section{Research Question}

Do chat embedding projections---measuring the cooperative, promise-making, homogeneity, and liar-like content of a player's messages---predict facial emotion during the subsequent contribution decision?

\section{Data and Timing}

Chat occurs \textbf{after} each round's contribution decision. Chat from round $N$ is paired with round $N+1$'s data. Facial emotion is measured via iMotions (webcam-based affect detection) during the \textbf{Contribute page} of round $N+1$.

\medskip
\noindent Timeline for a single player:
\begin{enumerate}[nosep]
    \item Round $N$: Contribute $\rightarrow$ Results $\rightarrow$ \textbf{Chat}
    \item Round $N+1$: \textbf{Facial emotion measured} $\rightarrow$ Contribute $\rightarrow$ Results $\rightarrow$ Chat
\end{enumerate}

Round 1 has no prior chat and is excluded, yielding $N = 1{,}791$ complete-case observations across 10 sessions and 5 supergames.

\section{Projection Variables}

Each player's chat messages are embedded (text-embedding-3-small), then projected onto direction vectors defined by group centroids:

\begin{itemize}[nosep]
    \item \textbf{Cooperative}: centroid of high-contributor chat $-$ low-contributor chat
    \item \textbf{Promise}: centroid of promise-making chat $-$ non-promise chat
    \item \textbf{Homogeneity}: within-group embedding similarity
    \item \textbf{Round-liar}: centroid of liar chat $-$ non-liar chat (liar $=$ promised $\wedge$ contributed $< 20$)
\end{itemize}

Multiple messages per player-round are averaged to a single projection score.

\section{Specification}

Each cell in the results table is a separate univariate regression:
\begin{equation*}
    \text{Emotion}_Y \;=\; \beta \cdot \text{Projection}_X \;+\; \gamma_1 \cdot \text{WordCount} \;+\; \gamma_2 \cdot \text{Sentiment} \;+\; \alpha_i \;+\; \delta_s \;+\; \varepsilon
\end{equation*}

\begin{itemize}[nosep]
    \item Fixed effects: player ($\alpha_i$), supergame segment ($\delta_s$)
    \item Standard errors clustered at session $\times$ segment $\times$ group level
    \item Controls: word count and VADER compound sentiment
\end{itemize}

\section{Results}

\begingroup
\centering
\begin{tabular}{lccccc}
\tabularnewline \midrule \midrule
& Valence & Joy & Anger & Contempt & Surprise \\
\midrule
Cooperative & 7.348$^{*}$ & 8.097$^{*}$ & -0.1019 & 0.1925 & -0.4536 \\
& (4.180) & (4.233) & (0.1461) & (1.025) & (0.7800) \\
\addlinespace
Promise & -12.70$^{***}$ & -12.37$^{***}$ & 0.0130 & 0.5090 & -0.5947 \\
& (3.690) & (3.312) & (0.1780) & (0.8120) & (0.5368) \\
\addlinespace
Homogeneity & 7.655$^{*}$ & 8.037$^{*}$ & -0.0427 & 0.5493 & -0.1848 \\
& (4.198) & (4.274) & (0.1468) & (1.007) & (0.8043) \\
\addlinespace
Round-liar & -16.16$^{**}$ & -13.42$^{**}$ & 0.0527 & 2.908$^{**}$ & -0.5271 \\
& (6.277) & (5.756) & (0.2423) & (1.305) & (0.8479) \\
\addlinespace
\midrule
Observations & 1,791 & 1,791 & 1,791 & 1,791 & 1,791 \\
Controls & \multicolumn{5}{c}{Word Count, Sentiment (compound)} \\
Fixed Effects & \multicolumn{5}{c}{Player, Segment} \\
Clustering & \multicolumn{5}{c}{Session $\times$ Segment $\times$ Group} \\
\midrule \midrule
\multicolumn{6}{l}{\emph{Standard errors in parentheses}} \\
\multicolumn{6}{l}{\emph{Signif. Codes: ***: 0.01, **: 0.05, *: 0.1}} \\
\end{tabular}
\par\endgroup


\medskip
Key findings:
\begin{itemize}[nosep]
    \item \textbf{Promise} projection significantly predicts lower Valence ($-12.70^{***}$) and Joy ($-12.37^{***}$)
    \item \textbf{Round-liar} projection predicts lower Valence ($-16.16^{**}$), lower Joy ($-13.42^{**}$), and higher Contempt ($2.908^{**}$)
    \item \textbf{Cooperative} and \textbf{Homogeneity} show weak positive effects on Valence/Joy (marginal significance)
    \item No projection significantly predicts Anger or Surprise
\end{itemize}

\section{Interpretation Puzzle: Round-Liar Projection}

The negative Round-liar $\rightarrow$ Valence result is initially surprising: why would liar-like language predict negative facial affect?

\subsection{Ruling out the obvious explanation}

A separate regression of a binary liar flag on valence during the \textbf{Chat page} (same round as the lie) found a \textbf{positive} effect---liars show higher valence while chatting, not lower. So actual liars are not driving the negative valence result in the next round.

\subsection{Ruling out social contagion}

We tested whether being lied to in round $N-1$ causes a player to adopt liar-like language in round $N$. Among non-liars whose groupmates lied in the previous round:
\begin{itemize}[nosep]
    \item Correlation with own projection: $r = 0.03$ (negligible)
    \item Mean projection: 0.065 (others lied) vs.\ 0.057 (others did not) --- minimal difference
\end{itemize}
Social contagion of liar-like language is not supported.

\subsection{The actual driver: non-liars with liar-like language}

Decomposing by liar status reveals the effect comes from \textbf{non-liars} (n $= 1{,}720$):

\begin{center}
\begin{tabular}{lcccc}
\toprule
& Mean Projection & Mean Valence & Mean Contribution & N \\
\midrule
Non-liars, high projection ($\geq$ p75) & 0.130 & 4.01 & 23.7 & 430 \\
Non-liars, low projection ($\leq$ p25) & 0.001 & 6.65 & 23.4 & 430 \\
\bottomrule
\end{tabular}
\end{center}

These two groups are behaviorally identical (same contributions, same cooperative/noncooperative split, similar sentiment). They differ only in \textbf{linguistic style}: high-projection non-liars use shorter messages (6.9 vs.\ 8.3 words) that happen to resemble liar language.

\subsection{Open question}

The round-liar direction vector, while defined by the liar/non-liar centroid difference, appears to capture a linguistic dimension---possibly terse or disengaged communication---that correlates with negative facial affect but is orthogonal to actual lying behavior. The projection is a weak predictor of actual lying ($r = 0.28$) and a weaker predictor of facial emotion ($r = -0.06$), but the regression coefficient is statistically significant with player fixed effects.

What this linguistic dimension represents substantively, and whether the round-liar direction is the right way to capture it, remains an open question.

\end{document}
